Modern GPU domain-specific languages (DSLs), such as Triton and TileLang, are increasingly used to implement specialized deep-learning kernels and as target languages for automated kernel-generation systems. Existing DSL-kernel evaluations commonly establish correctness through reference-based numerical validation. While this criterion is necessary, it does not characterize replacement quality, since a functionally valid kernel may still fall far below the latency and throughput of the optimized library operator it is intended to replace. 

We study this correctness–performance gap as an evaluation problem for GPU DSL kernels. Using 22 Triton and TileLang kernels from five operator categories on NVIDIA A100 and GH200 GPUs, we ask whether current correctness-based evaluation identifies kernels that are unsuitable as library replacements, why such failures occur, and how they can be detected without exhaustive benchmark coverage. The study yields three results. {First}, correctness-based evaluation can admit severe slowdowns: an idiomatic TileLang LayerNorm kernel passes KernelBench’s correctness check while running more than 300× slower than the PyTorch baseline. {Second}, the causes differ by kernel family. TileLang normalization and reduction slowdowns are mainly repairable authoring defects, such as sequential reductions and unnecessary dtype conversions, whereas convolution and large GEMM retain residual gaps after optimization due to code generation, autotuning coverage, and vendor-library algorithm selection. {Third}, library-relative efficiency and roofline utilization provide complementary screening criteria for identifying functionally valid but inefficient DSL kernels.